\ifdefined\gkver\else\def\gkver{student}\fi
\documentclass[\gkver]{gaokaozhenti}
\begin{document}

\chapter{2023年全国乙卷}
%% paperType: 理综物理部分

\begin{enumerate}
\item
%% number: 1
%% paperName: 2023年全国高考乙卷第 1 题 6 分
%% typeId: 1
%% type: 单选题
%% chapter: 运动和力
%% point: 变加速直线运动
%% method: 无
%% score: 6
%% degree: 550
%% duplicateId: 0
%% body:
一同学将排球自 O 点垫起，排球竖直向上运动，随后下落回到 O 点。设排球在运动过程中所受空气阻力大小和速度大小成正比。则该排球\xzanswer{B}
\fourchoices[answer=B]
{上升时间等于下落时间}
{被垫起后瞬间的速度最大}
{达到最高点时加速度为零}
{下落过程中做匀加速运动}
%% answer: B
%% memo:
\memoanswer{
A．上升过程和下降过程的位移大小相同，上升过程的末状态和下降过程的初状态速度均为零。对排球受力分析，上升过程的重力和阻力方向相同，下降过程中重力和阻力方向相反，根据牛顿第二定律可知，上升过程中任意位置的加速度比下降过程中对应位置的加速度大，则上升过程的平均加速度较大。由位移与时间关系可知，上升时间比下落时间短，A 错误；

B．上升过程排球做减速运动，下降过程排球做加速运动。在整个过程中空气阻力一直做负功，小球机械能一直在减小，下降过程中的最低点的速度小于上升过程的最低点的速度，故排球被垫起时的速度最大，B 正确；

C．达到最高点速度为零，空气阻力为零，此刻排球重力提供加速度不为零，C 错误；

D．下落过程中，排球速度在变，所受空气阻力在变，故排球所受的合外力在变化，排球在下落过程中做变加速运动，D 错误。

故选 B。
}

\item
%% number: 2
%% paperName: 2023年全国高考乙卷第 2 题 6 分
%% typeId: 1
%% type: 单选题
%% chapter: 曲线运动
%% point: 曲线运动条件
%% method: 无
%% score: 6
%% degree: 650
%% duplicateId: 0
%% body:
小车在水平地面上沿轨道从左向右运动，动能一直增加。如果用带箭头的线段表示小车在轨道上相应位置处所受合力，下列四幅图可能正确的是\xzanswer{D}
\onepicture{\includesvg{figs/02}}
%% answer: D
%% memo:
\memoanswer{
AB．小车做曲线运动，所受合外力指向曲线的凹侧，故 AB 错误；

CD．小车沿轨道从左向右运动，动能一直增加，故合外力与运动方向夹角为锐角，C 错误，D 正确。

故选 D。
}

\item
%% number: 3
%% paperName: 2023年全国高考乙卷第 3 题 6 分
%% typeId: 1
%% type: 单选题
%% chapter: 相对论简介
%% point: 质速关系和质能关系
%% method: 无
%% score: 6
%% degree: 650
%% duplicateId: 0
%% body:
2022 年 10 月，全球众多天文设施观测到迄今最亮伽马射线暴，其中我国的“慧眼”卫星、“极目”空间望远镜等装置在该事件观测中作出重要贡献。由观测结果推断，该伽马射线暴在 1 分钟内释放的能量量级为 10$^{48}$ J。假设释放的能量来自于物质质量的减少，则每秒钟平均减少的质量量级为（光速为 3$\times$10$^{8}$ $\Ums$）\xzanswer{C}
\fourchoices[answer=C]
{10$^{19}$ $\Ukg$}
{10$^{24}$ $\Ukg$}
{10$^{29}$ $\Ukg$}
{10$^{34}$ $\Ukg$}
%% answer: C
%% memo:
\memoanswer{
根据质能方程  $E = mc^{2}$ 可知，则每秒钟平均减少的质量为

$$ \Delta m=\frac{E_{0}}{60c^{2}}=\frac{10^{48}}{60\times\left(3\times10^{8}\right)^{2}}=\frac{10^{30}}{5.4}kg $$

则每秒钟平均减少的质量量级为  $10^{29}$ $\Ukg$。

故选 C。
}

\item
%% number: 4
%% paperName: 2023年全国高考乙卷第 4 题 6 分
%% typeId: 1
%% type: 单选题
%% chapter: 电磁感应
%% point: 电磁感应中图象
%% method: 无
%% score: 6
%% degree: 550
%% duplicateId: 0
%% body:
一学生小组在探究电磁感应现象时，进行了如下比较实验。用图（a）所示的缠绕方式，将漆包线分别绕在几何尺寸相同的有机玻璃管和金属铝管上，漆包线的两端与电流传感器接通。两管皆竖直放置，将一很小的强磁体分别从管的上端由静止释放，在管内下落至管的下端。实验中电流传感器测得的两管上流过漆包线的电流 $I$ 随时间 $t$ 的变化分别如图（b）和图（c）所示，分析可知\xzanswer{A}
\twopicture[labelA=2023全国乙卷04a, labelB=2023全国乙卷04b]
{figs/04a.png}
{figs/04b.png}
\fourchoices[answer=A]
{图（c）是用玻璃管获得的图像}
{在铝管中下落，小磁体做匀变速运动}
{在玻璃管中下落，小磁体受到的电磁阻力始终保持不变}
{用铝管时测得的电流第一个峰到最后一个峰的时间间隔比用玻璃管时的短}
%% answer: A
%% memo:
\memoanswer{
A．强磁体在铝管中运动，铝管会形成涡流，玻璃是绝缘体故强磁体在玻璃管中运动，玻璃管不会形成涡流。强磁体在铝管中加速后很快达到平衡状态，做匀速直线运动，而玻璃管中的磁体则一直做加速运动，故由图像可知图（c）的脉冲电流峰值不断增大，说明强磁体的速度在增大，与玻璃管中磁体的运动情况相符，A 正确；

B．在铝管中下落，脉冲电流的峰值一样，磁通量的变化率相同，故小磁体做匀速运动，B 错误；

C．在玻璃管中下落，玻璃管为绝缘体，线圈的脉冲电流峰值增大，电流不断在变化，故小磁体受到的电磁阻力在不断变化，C 错误；

D．强磁体分别从管的上端由静止释放，在铝管中，磁体在线圈间做匀速运动，玻璃管中磁体在线圈间做加速运动，故用铝管时测得的电流第一个峰到最后一个峰的时间间隔比用玻璃管时的长，D 错误。

故选 A。
}

\item
%% number: 5
%% paperName: 2023年全国高考乙卷第 5 题 6 分
%% typeId: 1
%% type: 单选题
%% chapter: 磁场
%% point: 带电粒子在复合场中的运动
%% method: 无
%% score: 6
%% degree: 550
%% duplicateId: 0
%% body:
如图，一磁感应强度大小为 $B$ 的匀强磁场，方向垂直于纸面（$xOy$ 平面）向里，磁场右边界与 $x$ 轴垂直。一带电粒子由 $O$ 点沿 $x$ 正向入射到磁场中，在磁场另一侧的 S 点射出，粒子离开磁场后，沿直线运动打在垂直于 $x$ 轴的接收屏上的 P 点；SP = $l$，S 与屏的距离为 $\frac{l}{2}$，与 $x$ 轴的距离为 $a$。如果保持所有条件不变，在磁场区域再加上电场强度大小为 $E$ 的匀强电场，该粒子入射后则会沿 $x$ 轴到达接收屏。该粒子的比荷为\xzanswer{A}
\onepicture{figs/05.png}
\fourchoices[answer=A]
{$\frac{E}{{2a{B^2}}}$}
{$\frac{E}{{a{B^2}}}$}
{$\frac{B}{{2a{E^2}}}$}
{$\frac{B}{{a{E^2}}}$}
%% answer: A
%% memo:
\memoanswer{
由题知，一带电粒子由 O 点沿 x 正向入射到磁场中，在磁场另一侧的 S 点射出，则根据几何关系可知粒子做圆周运动的半径  $r = 2a$

则粒子做圆周运动有  $qvB = m \frac{v^{2}}{r}$

则有  $\frac{q}{m} = \frac{v}{2a \cdot B}$

如果保持所有条件不变，在磁场区域再加上电场强度大小为 E 的匀强电场，该粒子入射后则会沿 x 轴到达接收屏，则有

 $Eq = qvB$

联立有

 $\frac{q}{m} = \frac{E}{2a \cdot B^{2}}$

故选 A。
}

\item
%% number: 6
%% paperName: 2023年全国高考乙卷第 6 题 6 分
%% typeId: 2
%% type: 多选题
%% chapter: 电场
%% point: 带电粒子的轨迹类问题
%% method: 无
%% score: 6
%% degree: 550
%% duplicateId: 0
%% body:
在 O 点处固定一个正点电荷，P 点在 O 点右上方。从 P 点由静止释放一个带负电的小球，小球仅在重力和该点电荷电场力作用下在竖直面内运动，其一段轨迹如图所示。M、N 是轨迹上的两点，OP > OM，OM = ON，则小球\xzanswer{BC}
\onepicture{figs/06.png}
\fourchoices[answer=BC]
{在运动过程中，电势能先增加后减少}
{在 P 点的电势能大于在 N 点的电势能}
{在 M 点的机械能等于在 N 点的机械能}
{从 M 点运动到 N 点的过程中，电场力始终不做功}
%% answer: BC
%% memo:
\memoanswer{
ABC．由题知，OP > OM，OM = ON，则根据点电荷的电势分布情况可知

$\varphi_{M} = \varphi_{N} > \varphi_{P}$

则带负电的小球在运动过程中，电势能先减小后增大，且

$E_{pP} > E_{pM} = E_{pN}$

则带负电的小球在 M 点的机械能等于在 N 点的机械能，A 错误、BC 正确；

D．从 M 点运动到 N 点的过程中，电场力先做正功后做负功，D 错误。

故选 BC。
}

\item
%% number: 7
%% paperName: 2023年全国高考乙卷第 7 题 6 分
%% typeId: 2
%% type: 多选题
%% chapter: 电路
%% point: 串并联电路
%% method: 无
%% score: 6
%% degree: 450
%% duplicateId: 0
%% body:
黑箱外有编号为 1、2、3、4 的四个接线柱，接线柱 1 和 2、2 和 3、3 和 4 之间各接有一个电阻，在接线柱间还接有另外一个电阻 $R$ 和一个直流电源。测得接线柱之间的电压 $U_{12}$ = 3.0 $\UV$，$U_{23}$ = 2.5 $\UV$，$U_{34} = -$ 1.5 $\UV$。符合上述测量结果的可能接法是\xzanswer{CD}
\onepicture{\includesvg{figs/07}}
\fourchoices[answer=CD]
{电源接在 1、4 之间，$R$ 接在 1、3 之间}
{电源接在 1、4 之间，$R$ 接在 2、4 之间}
{电源接在 1、3 之间，$R$ 接在 1、4 之间}
{电源接在 1、3 之间，$R$ 接在 2、4 之间}
%% answer: CD
%% memo:
\memoanswer{
A．根据题意画出电路图，如下

\onepicture[w=6cm]{figs/07_an1.jpg}

可见  $U_{34} > 0$，A 错误；

B．根据题意画出电路图，如下

\onepicture[w=6cm]{figs/07_an2.jpg}

可见  $U_{34} > 0$，B 错误；

C. 根据题意画出电路图，如下

\onepicture[w=6cm]{figs/07_an3.jpg}

可见上述接法可符合上述测量结果，C 正确；

D．根据题意画出电路图，如下

\onepicture[w=6cm]{figs/07_an4.jpg}

可见上述接法可符合上述测量结果，D 正确。

故选 CD。
}

\item
%% number: 8
%% paperName: 2023年全国高考乙卷第 8 题 6 分
%% typeId: 2
%% type: 多选题
%% chapter: 功和能
%% point: 动能定理
%% method: 无
%% score: 6
%% degree: 450
%% duplicateId: 0
%% body:
如图，一质量为 $M$、长为 $l$ 的木板静止在光滑水平桌面上，另一质量为 $m$ 的小物块（可视为质点）从木板上的左端以速度 $v_{0}$ 开始运动。已知物块与木板间的滑动摩擦力大小为 $f$，当物块从木板右端离开时\xzanswer{BD}
\onepicture{figs/08.png}
\fourchoices[answer=BD]
{木板的动能一定等于 $fl$}
{木板的动能一定小于 $fl$}
{物块的动能一定大于 $\frac{1}{2}mv_{0}^{2} - fl$}
{物块的动能一定小于 $\frac{1}{2}mv_{0}^{2} - fl$}
%% answer: BD
%% memo:
\memoanswer{
当物块从木板右端离开时，对 m 有

 $-fx_{m}=\frac{1}{2}mv_{m}^{2}-\frac{1}{2}mv_{0}^{2}$

对 M 有

 $fx_{M}=\frac{1}{2}mv_{M}^{2}$

其中  $l = x_{m} - x_{M}$

AB．由于  $l > x_{M}$，则根据以上分析可知木板的动能一定小于  $fl$，A 错误、B 正确；

CD．根据以上分析，联立有

 $\frac{1}{2}mv_{0}^{2}-fl=\frac{1}{2}mv_{m}^{2}+\frac{1}{2}mv_{M}^{2}$

则物块的动能一定小于  $\frac{1}{2}mv_{0}^{2}-fl$，C 错误、D 正确。

故选 BD。


（一）必考题：共129分。
}

\item
%% number: 9
%% paperName: 2023年全国高考乙卷第 9 题 5 分
%% typeId: 4
%% type: 实验题
%% chapter: 力物体的平衡
%% point: 验证平行四边形实验
%% method: 无
%% score: 5
%% degree: 750
%% duplicateId: 0
%% body:
在“验证力的平行四边形定则”的实验中使用的器材有：木板、白纸、两个标准弹簧测力计、橡皮条、轻质小圆环、刻度尺、铅笔、细线和图钉若干。完成下列实验步骤：

\begin{steps}
	\item
	用图钉将白纸固定在水平木板上。
	\item
	将橡皮条的一端固定在木板上，另一端系在轻质小圆环上。将两细线也系在小圆环上，它们的另一端均挂上测力计。用互成一定角度、方向平行于木板、大小适当的力拉动两个测力计，小圆环停止时由两个测力计的示数得到两拉力 $F_{1}$ 和 $F_{2}$ 的大小，并 \tkanswer{CD} 。（多选，填正确答案标号）
	\item
	撤掉一个测力计，用另一个测力计把小圆环拉到 \tkanswer{相同位置} ，由测力计的示数得到拉力 $F$ 的大小，沿细线标记此时 $F$ 的方向。
	\item
	选择合适标度，由步骤②的结果在白纸上根据力的平行四边形定则作 $F_{1}$ 和 $F_{2}$ 的合成图，得出合力 $F'$ 的大小和方向；按同一标度在白纸上画出力 $F$ 的图示。
	\item
	比较 $F'$ 和 $F$ 的 \tkanswer{大小和方向} ，从而判断本次实验是否验证了力的平行四边形定则。
\end{steps}
%% answer: 
\jdanswer{
CD，相同位置，大小和方向
}
%% memo:
\memoanswer{
① 用图钉将白纸固定在水平木板上，以便在白纸上记录各力的作用点与方向。

② 将橡皮条的一端固定在木板上，另一端系在轻质小圆环上。将两细线也系在小圆环上，它们的另一端均挂上测力计。用互成一定角度、方向平行于木板、大小适当的力拉动两个测力计，小圆环停止时由两个测力计的示数得到两拉力 $F_{1}$ 和 $F_{2}$ 的大小，还需要用铅笔在白纸上标记出小圆环的位置以及用铅笔在白纸上标记出两细线的方向。故选 CD。

③ 撤掉一个测力计，用另一个测力计把小圆环拉到相同位置，由测力计的示数得到拉力 $F$ 的大小，沿细线标记此时 $F$ 的方向。

④ 选择合适标度，由步骤②的结果在白纸上根据力的平行四边形定则作 $F_{1}$ 和 $F_{2}$ 的合成图，得出合力 $F'$ 的大小和方向；按同一标度在白纸上画出力 $F$ 的图示。

⑤ 比较 $F'$ 和 $F$ 的大小和方向，从而判断本次实验是否验证了力的平行四边形定则。
}

\item
%% number: 10
%% paperName: 2023年全国高考乙卷第 10 题 10 分
%% typeId: 4
%% type: 实验题
%% chapter: 电路
%% point: 测量电阻率
%% method: 无
%% score: 10
%% degree: 550
%% duplicateId: 0
%% body:
一学生小组测量某金属丝（阻值约十几欧姆）的电阻率。现有实验器材：螺旋测微器、米尺、电源 $E$、电压表（内阻非常大）、定值电阻 $R_{0}$（阻值 10.0 $\UO$）、滑动变阻器 $R$、待测金属丝、单刀双掷开关 K、开关 S、导线若干。图（a）是学生设计的实验电路原理图。完成下列填空：

\begin{enumerate}
	\item
	实验时，先将滑动变阻器 $R$ 接入电路的电阻调至最大，闭合 S。
	\item
	将 K 与 1 端相连，适当减小滑动变阻器 $R$ 接入电路的电阻，此时电压表读数记为 $U_{1}$，然后将 K 与 2 端相连，此时电压表读数记为 $U_{2}$。由此得到流过待测金属丝的电流$I$ = \_\_\_\_\_\_\_，金属丝的电阻 $r$ = \_\_\_\_\_。（结果均用 $R_{0}$、$U_{1}$、$U_{2}$ 表示）
	\item
	继续微调 $R$，重复（2）的测量过程，得到多组测量数据，如下表所示：

	\begin{center}
	\begin{tblr}{|c|c|c|c|c|c|}
	\hline
	$U_{1}$（$\UmV$） & 0.57 & 0.71 & 0.85 & 1.14 & 1.43 \\
	\hline
	$U_{2}$（$\UmV$） & 0.97 & 1.21 & 1.45 & 1.94 & 2.43 \\
	\hline
	\end{tblr}
	\end{center}
	\item
	利用上述数据，得到金属丝的电阻 $r$ = 14.2 $\UO$。
	\item
	用米尺测得金属丝长度 $L$ = 50.00 $\Ucm$。用螺旋测微器测量金属丝不同位置的直径，某次测量的示数如图（b）所示，该读数为 $d$ = \_\_\_\_\_\_ $\Umm$。多次测量后，得到直径的平均值恰与 $d$ 相等。
	\item
	由以上数据可得，待测金属丝所用材料的电阻率 $\rho$ = \_\_\_\_\_\_$\times10^{-7}\UOm$。（保留 2 位有效数字）
\end{enumerate}
\onepicture[align=right]{figs/10.png}
%% answer: 
\jdanswer{
\begin{enumerate}
	\item
	$\frac{{{U_2} - {U_1}}}{{{R_0}}}$，$\frac{{{U_1}{R_0}}}{{{U_2} - {U_1}}}$

	\item
	0.150

	\item
	5.0

\end{enumerate}
}
%% memo:
\memoanswer{
（2）根据题意可知， $R_{0}$ 两端的电压为 $U = U_{2} - U_{1}$

则流过 $R_{0}$ 即流过待测金属丝的电流 $I = \frac{U}{R_0} = \frac{U_2 - U_1}{R_0}$

金属丝的电阻 $r = \frac{U_{1}}{I}$

联立可得 $r = \frac{U_{1}R_{0}}{U_{2} - U_{1}}$

（5）螺旋测微器的读数为 $d = 15.0 \times 0.01 \mathrm{~mm} = 0.150 \mathrm{~mm}$

（6）根据电阻定律 $r = \rho \frac{L}{S}$

又 $S = \pi \cdot (\frac{d}{2})^2$

代入数据联立解得 $\rho = 5.0 \times 10^{-7}\UOm$
}

\item
%% number: 11
%% paperName: 2023年全国高考乙卷第 11 题 12 分
%% typeId: 6
%% type: 计算题
%% chapter: 电场
%% point: 电场叠加
%% method: 无
%% score: 12
%% degree: 550
%% duplicateId: 0
%% body:
如图，等边三角形 $\triangle$ABC 位于竖直平面内，AB 边水平，顶点 C 在 AB 边上方，3 个点电荷分别固定在三角形的三个顶点上。已知 AB 边中点 M 处的电场强度方向竖直向下，BC 边中点 N 处的电场强度方向竖直向上，A 点处点电荷的电荷量的绝对值为 $q$，求
\begin{enumerate}
	\item
	B 点处点电荷的电荷量的绝对值并判断 3 个点电荷的正负；
	\item
	C 点处点电荷的电荷量。
\end{enumerate}
\onepicture[align=right]{figs/11.png}
%% answer: 
\jdanswer{
\begin{enumerate}
	\item
	$q$，A、B、C 均为正电荷

	\item
	$\frac{{3 - \sqrt 3 }}{3}q$

\end{enumerate}
}
%% memo:
\memoanswer{
\onepicture[w=6cm]{figs/11_an.png}



\onepicture[w=6cm]{figs/11_an2.png}
}

\item
%% number: 12
%% paperName: 2023年全国高考乙卷第 12 题 20 分
%% typeId: 6
%% type: 计算题
%% chapter: 运动和力
%% point: 动量和能量
%% method: 无
%% score: 20
%% degree: 350
%% duplicateId: 0
%% body:
如图，一竖直固定的长直圆管内有一质量为 $M$ 的静止薄圆盘，圆盘与管的上端口距离为 $l$，圆管长度为 20$l$。一质量为 $m = \frac{1}{3}M$ 的小球从管的上端口由静止下落，并撞在圆盘中心，圆盘向下滑动，所受滑动摩擦力与其所受重力大小相等。小球在管内运动时与管壁不接触，圆盘始终水平，小球与圆盘发生的碰撞均为弹性碰撞且碰撞时间极短。不计空气阻力，重力加速度大小为 $g$。求
\begin{enumerate}
	\item
	第一次碰撞后瞬间小球和圆盘的速度大小；
	\item
	在第一次碰撞到第二次碰撞之间，小球与圆盘间的最远距离；
	\item
	圆盘在管内运动过程中，小球与圆盘碰撞的次数。
\end{enumerate}
\onepicture[align=right]{figs/12.png}
%% answer: 
\jdanswer{
\begin{enumerate}
	\item
	小球速度大小 $\frac{{\sqrt {2gl} }}{2}$，圆盘速度大小 $\frac{{\sqrt {2gl} }}{2}$

	\item
	$l$

	\item
	4

\end{enumerate}
}
%% memo:
\memoanswer{
（1）过程 1：小球释放后自由下落，下降 l，根据机械能守恒定律

 $mgl = \frac{1}{2}mv_{0}^{2}$

解得  $v_{0} = 2gl$

过程 2：小球以  $2gl$ 与静止圆盘发生弹性碰撞，根据能量守恒定律和动量守恒定律分别有

 $\frac{1}{2}mv_{0}^{2} = \frac{1}{2}mv_{1}^{2} + \frac{1}{2}Mv_{1}^{2}$

 $mv_{0} = mv_{1} + Mv_{1}^{'}$

解得

$$ v_{1}=\frac{m-M}{m+M}v_{0}=-\frac{\sqrt{2gl}}{2} $$

$$ v_{1}^{\prime}=\frac{1}{2}v_{0}=\frac{\sqrt{2gl}}{2} $$

即小球碰后速度大小  $\frac{2gl}{2}$，方向竖直向上，圆盘速度大小为  $\frac{2gl}{2}$，方向竖直向下；

（2）第一次碰后，小球做竖直上抛运动，圆盘摩擦力与重力平衡，匀速下滑，所以只要圆盘下降速度比小球快，二者间距就不断增大，当二者速度相同时，间距最大，即

 $v_{1} + gt = v_{1}^{'}$

解得

$$ t=\frac{v_{1}^{\prime}-v_{1}}{g}=\frac{v_{0}}{g} $$

根据运动学公式得最大距离为

$$ d_{\max}=x_{ \text{盘} }-x_{ \text{球} }=v_{1}^{\prime}t-\left(v_{1}t-\frac{1}{2}gt^{2}\right)=\frac{v_{0}^{2}}{2g}=l $$

（3）第一次碰撞后到第二次碰撞时，两者位移相等，则有

 $x_{\text{盘}1}=x_{\text{球}1}$

即

 $v_{1}t_{1}+\frac{1}{2}gt_{1}^{2}=v_{1}^{'}t_{1}$

解得

 $t_{1}=\frac{2v_{0}}{g}$

此时小球的速度

 $v_{2}=v_{1}+gt_{1}=\frac{3}{2}v_{0}$

圆盘的速度仍为  $v_{1}^{'}$，这段时间内圆盘下降的位移

 $x_{\text{盘}1}=v_{1}^{'}t_{1}=\frac{v_{0}^{2}}{g}=2l$

之后第二次发生弹性碰撞，根据动量守恒

 $mv_{2}+Mv_{1}^{'}=mv_{2}^{'}+Mv_{2}^{''}$

根据能量守恒

 $\frac{1}{2}mv_{2}^{2}+\frac{1}{2}Mv_{1}'^{2}=\frac{1}{2}mv_{2}'^{2}+\frac{1}{2}Mv_{2}''^{2}$

联立解得

$$ \dot{v_{2}}=0,\quad\dot{v_{2}}=v_{0} $$

同理可得当位移相等时

$$ x_{ \text{盘} 2}=x_{ \text{球} 2} $$

$$ v_{2}^{\prime \prime}t_{2}=\frac{1}{2}gt_{2}^{2} $$

解得

$$ t_{2}=\frac{2v_{0}}{g} $$

圆盘向下运动

$$ x_{ \text{盘} 2}=v_{2}^{\prime \prime}t_{2}=\frac{2v_{0}^{2}}{g}=4l $$

此时圆盘距下端关口 13l，之后二者第三次发生碰撞，碰前小球的速度

$$ v_{3}=gt_{2}=2v_{0} $$

有动量守恒

$$ m v_{3}+M v_{2}^{\prime \prime}=m v_{3}^{\prime}+M v_{3}^{\prime \prime} $$

机械能守恒

$$ \frac{1}{2}m v_{3}^{2}+\frac{1}{2}M v_{2}^{\prime \prime2}=\frac{1}{2}m v_{3}^{\prime2}+\frac{1}{2}M v_{3}^{\prime \prime2} $$

得碰后小球速度为

$$ v_{3}^{\prime}=\frac{v_{0}}{2} $$

圆盘速度

$$ v_{3}^{\prime \prime}=\frac{3v_{0}}{2} $$

当二者即将四次碰撞时  $x_{\text{盘}3}=x_{\text{球}3}$

即

$$ v_{3}^{\prime \prime}t_{3}=v_{3}^{\prime}t_{3}+\frac{1}{2}gt_{3}^{2} $$

得

$$ t_{3}=\frac{2v_{0}}{g}=t_{1}=t_{2} $$

在这段时间内，圆盘向下移动

$$ x_{ \text{盘} 3}=v_{3}^{\prime \prime}t_{3}=\frac{3v_{0}^{2}}{g}=6l $$

此时圆盘距离下端管口长度为  $20l - 1l - 2l - 4l - 6l = 7l$

此时可得出圆盘每次碰后到下一次碰前，下降距离逐次增加 2l，故若发生下一次碰撞，圆盘将向下移动  $x_{\text{盘}4}=8l$

则第四次碰撞后落出管口外，因此圆盘在管内运动的过程中，小球与圆盘的碰撞次数为 4 次。

（二）选考题：共 45 分。请考生从给出的 2 道物理题、2 道化学题、2 道生物题中每科任选一题作答，并用 2B 铅笔在答题卡上把所选题目题号后的方框涂黑。注意所做题目的题号必须与所涂题目的题号一致，并且在解答过程中写清每问的小题号，在答题卡指定位置答题。如果多做则每学科按所做的第一题计分。
}

\item
%% number: 13
%% paperName: 2023年全国高考乙卷第 13 题 5 分
%% typeId: 2
%% type: 多选题
%% chapter: 气体性质
%% point: 物体的内能
%% method: 无
%% score: 5
%% degree: 650
%% duplicateId: 0
%% body:
对于一定量的理想气体，经过下列过程，其初始状态的内能与末状态的内能可能相等的是\xzanswer{ACD}
\fivechoices[answer=ACD]
{等温增压后再等温膨胀}
{等压膨胀后再等温压缩}
{等容减压后再等压膨胀}
{等容增压后再等压压缩}
{等容增压后再等温膨胀}
%% answer: ACD
%% memo:
\memoanswer{
A．对于一定质量的理想气体内能由温度决定，故等温增压和等温膨胀过程温度均保持不变，内能不变，故 A 正确；

B．根据理想气体状态方程可知：等压膨胀后气体温度升高，内能增大，等温压缩温度不变，内能不变，故末状态与初始状态相比内能增加，故 B 错误；

C．根据理想气体状态方程可知：等容减压过程温度降低，内能减小；等压膨胀过程温度升高，末状态的温度有可能和初状态的温度相等，内能相等，故 C 正确；

D．根据理想气体状态方程可知：等容增压过程温度升高；等压压缩过程温度降低，末状态的温度有可能和初状态的温度相等，内能相等，故 D 正确；

E．根据理想气体状态方程可知：等容增压过程温度升高；等温膨胀过程温度不变，故末状态的内能大于初状态的内能，故 E 错误。

故选 ACD。
}

\item
%% number: 13
%% paperName: 2023年全国高考乙卷第 13 题 0 分
%% typeId: 6
%% type: 计算题
%% chapter: 气体性质
%% point: 气体连接体
%% method: 无
%% score: 0
%% degree: 550
%% duplicateId: 0
%% body:
如图，竖直放置的封闭玻璃管由管径不同、长度均为 20 $\Ucm$ 的 A、B 两段细管组成，A 管的内径是 B 管的 2 倍，B 管在上方。管内空气被一段水银柱隔开。水银柱在两管中的长度均为 10 $\Ucm$。现将玻璃管倒置使 A 管在上方，平衡后，A 管内的空气柱长度改变 1 $\Ucm$。求 B 管在上方时，玻璃管内两部分气体的压强。（气体温度保持不变，以 cmHg 为压强单位）
\onepicture[align=right]{figs/13.png}
%% answer: 
\jdanswer{
$p_{A}$ = 135 cmHg，$p_{B}$ = 115 cmHg
}
%% memo:
\memoanswer{
A．对于一定质量的理想气体内能由温度决定，故等温增压和等温膨胀过程温度均保持不变，内能不变，故 A 正确；

B．根据理想气体状态方程可知：等压膨胀后气体温度升高，内能增大，等温压缩温度不变，内能不变，故末状态与初始状态相比内能增加，故 B 错误；

C．根据理想气体状态方程可知：等容减压过程温度降低，内能减小；等压膨胀过程温度升高，末状态的温度有可能和初状态的温度相等，内能相等，故 C 正确；

D．根据理想气体状态方程可知：等容增压过程温度升高；等压压缩过程温度降低，末状态的温度有可能和初状态的温度相等，内能相等，故 D 正确；

E．根据理想气体状态方程可知：等容增压过程温度升高；等温膨胀过程温度不变，故末状态的内能大于初状态的内能，故 E 错误。

故选 ACD。

（10 分）如图，竖直放置的封闭玻璃管由管径不同、长度均为 20 $\Ucm$ 的 A、B 两段细管组成，A 管的内径是 B 管的 2 倍，B 管在上方。管内空气被一段水银柱隔开。水银柱在两管中的长度均为 10 $\Ucm$。现将玻璃管倒置使 A 管在上方，平衡后，A 管内的空气柱长度改变 1 $\Ucm$。求 B 管在上方时，玻璃管内两部分气体的压强。（气体温度保持不变，以 cmHg 为压强单位）

【详解】设 B 管在上方时上部分气压为  $p_{B}$，则此时下方气压为  $p_{A}$，此时有  $p_{A}=p_{B}+20$

倒置后 A 管气体压强变小，即空气柱长度增加 1 $\Ucm$，A 管中水银柱减小 1 $\Ucm$，又因为  $S_{A} = 2S_{B}$

可知 B 管水银柱增加 2 $\Ucm$，空气柱减小 2 $\Ucm$；设此时两管的压强分别为  $p_{A}'$、 $p_{B}'$，所以有  $p_{A}' + 21 = p_{B}'$

倒置前后温度不变，根据玻意耳定律对 A 管有  $p_{A}S_{A}L_{A}=p_{A}^{\prime}S_{A}L_{A}^{\prime}$

对 B 管有  $p_{\mathrm{B}}S_{\mathrm{B}}L_{\mathrm{B}}=p_{\mathrm{B}}^{\prime}S_{\mathrm{B}}L_{\mathrm{B}}^{\prime}$

其中  $L_{A}^{\prime}=10\Ucm+1\Ucm=11\Ucm$， $L_{B}^{\prime}=10\Ucm-2\Ucm=8\Ucm$

联立以上各式解得： $p_{A}=135\ cmHg$， $p_{B}=115\ cmHg$
}

\item
%% number: 14
%% paperName: 2023年全国高考乙卷第 14 题 5 分
%% typeId: 2
%% type: 多选题
%% chapter: 机械波
%% point: 波与振动图象
%% method: 无
%% score: 5
%% degree: 550
%% duplicateId: 0
%% body:
一列简谐横波沿 $x$ 轴传播，图（a）是 $t$ = 0 时刻的波形图；P 是介质中位于 $x$ = 2 $\Um$ 处的质点，其振动图像如图（b）所示。下列说法正确的是\xzanswer{ABE}
\onepicture{\includesvg{figs/14}}
\fivechoices[answer=ABE]
{波速为 2 $\Ums$}
{波向左传播}
{波的振幅是 10 $\Ucm$}
{$x$ = 3 $\Um$ 处的质点在 $t$ = 7 $\Us$ 时位于平衡位置}
{质点 P 在 0 \textasciitilde{} 7 $\Us$ 时间内运动的路程为 70 $\Ucm$}
%% answer: ABE
%% memo:
\memoanswer{
A．由图 a 可知波长为 4 $\Um$，由图 b 可知波的周期为 2 $\Us$，则波速为  $v = \frac{\lambda}{T} = \frac{4}{2} \, \Ums = 2 \, \Ums$

故 A 正确；

B．由图乙可知 t=0 时，P 点向下运动，根据 “上下坡” 法可知波向左传播，故 B 正确；

C．由图（a）可知波的振幅为 5 $\Ucm$，故 C 错误；

DE．根据图（a）可知 t=0 时  $x=3\ m$ 处的质点位于波谷处，由于

 $t=7\Us=3T+\frac{1}{2}T$

可知在  $t=7\ s$ 时质点位于波峰处；质点 P 运动的路程为

 $s=3\times4A+\frac{1}{2}\times4A=70\Ucm$

故 D 错误，E 正确；

故选 ABE。

（10 分）如图，一折射率为 2 的棱镜的横截面为等腰直角三角形  $\triangle ABC$， $AB = AC = l$，BC 边所在底面上镀有一层反射膜。一细光束沿垂直于 BC 方向经 AB 边上的 M 点射入棱镜，若这束光被 BC 边反射后恰好射向顶点 A，求 M 点到 A 点的距离。

【详解】由题意可知做出光路图如图所示

光线垂直于 BC 方向射入，根据几何关系可知入射角为  $45^\circ$；由于棱镜折射率为 2，根据  $n = \frac{\sin i}{\sin r}$

有

 $\sin r = \frac{1}{2}$

则折射角为  $30^\circ$； $\angle BMO = 60^\circ$，因为  $\angle B = 45^\circ$，所以光在 BC 面的入射角为

 $\theta = 90^\circ - \left(180^\circ - 60^\circ - 45^\circ\right) = 15^\circ$

根据反射定律可知

 $\angle MOA = 2\theta = 30^\circ$

根据几何关系可知  $\angle BAO = 30^\circ$，即  $\Delta MAO$ 为等腰三角形，则

 $\frac{MO}{AO} = \frac{\sqrt{3}}{3}$

又因为  $\Delta BOM$ 与  $\Delta CAO$ 相似，故有

 $\frac{BM}{AC} = \frac{MO}{AO}$

由题知  $AB = AC = l$

联立可得  $BM = \frac{\sqrt{3}}{3} AC = \frac{\sqrt{3}}{3} l$

所以 M 到 A 点的距离为

 $x = MA = l - BM = \frac{3 - \sqrt{3}}{3} l$
}

\item
%% number: 14
%% paperName: 2023年全国高考乙卷第 14 题 10 分
%% typeId: 6
%% type: 计算题
%% chapter: 光学
%% point: 光的折射
%% method: 无
%% score: 10
%% degree: 550
%% duplicateId: 0
%% body:
如图，一折射率为 $\sqrt 2 $ 的棱镜的横截面为等腰直角三角形 $\triangle$ABC，AB = AC = $l$，BC 边所在底面上镀有一层反射膜。一细光束沿垂直于 BC 方向经 AB 边上的 M 点射入棱镜，若这束光被 BC 边反射后恰好射向顶点 A，求 M 点到 A 点的距离。
\onepicture[align=right]{figs/14.png}
%% answer: 
\jdanswer{
$\frac{{3 - \sqrt 3 }}{3}l$
}
%% memo:
\memoanswer{
\onepicture[w=6cm]{figs/14_an.png}
}

\end{enumerate}

\end{document}
